\chapter{社会与人造智能的计算动力学：制度网络、交通流与语言模型}
\label{chp:social_artificial_intelligence}

生物智能把感知、调节和结构适应实现于有机体之内，社会与人造系统则把相同类型的计算问题分散到制度、基础设施、通信协议、模型参数和外部执行者之间。本章不把市场、城市交通或大语言模型拟人化，也不因为它们呈现宏观秩序就断言存在统一主体。我们关心的是更可检验的问题：哪些对象构成状态，哪些接口传递信息，局部更新如何形成全局分布，目标和约束由谁设置，结构何时改变，以及观测到的稳定、效率或预测能力究竟能支持什么结论。为此，我们以计算动力学统一描述连续量、图结构、离散事件和慢速制度变量，同时严格区分局部结构网络与全局分布表示。HSF--HD在此提供一般状态动力学的建模语言；AgentOS只作为后文可替换的工程实例，不能反向规定所有社会系统和人工系统的唯一机制。

\section{市场作为制度约束下的分布式估计与资源配置系统}
\label{sec:section_4093}

我们先定义市场的计算对象。设主体集合为 $V_t=\{1,\ldots,N_t\}$，可交易资产集合为 $A_t=\{1,\ldots,m_t\}$，二者都允许随时间进入或退出。主体 $i$ 在离散决策时刻 $t_k$ 的局部状态写成
\[
x_i^k=(b_i^k,q_i^k,\hat v_i^k,r_i^k,c_i^k),
\]
其中 $b_i^k\in\mathbb{R}$ 是以计价单位表示的现金头寸，$q_i^k\in\mathbb{R}^{m_k}$ 是资产数量，$\hat v_i^k\in\mathbb{R}^{m_k}$ 是条件估值，$r_i^k$ 是风险与期限状态，$c_i^k$ 是可用信息和合规权限。资产的法律权利、清偿顺序、抵押关系和交易资格构成类型化契约图 $G_k^{\mathrm{con}}=(V_k^{\mathrm{con}},E_k^{\mathrm{con}})$；订单、成交和资金划拨则是带时间戳的事件流，不与图节点混作同一种变量。价格 $p^k\in\mathbb{R}_{+}^{m_k}$、盘口深度、成交量和信用利差是从事件流聚合得到的全局分布表示，而企业控制权、债权优先级和支付通道属于局部结构网络。两者通过资产身份、账户权限、清算规则和读出窗口联合：一笔价格变化可以更新估值，却不会自动改写所有权；一次违约既改变现金流分布，也可能触发契约图上的边删除和优先级执行。

在这一类型划分下，资产不需要被解释为“形与质的张量积”，货币也不是为方程提供动力的物理能量。资产是可验证权利与条件现金流的组合，信用是跨期承诺和违约风险的制度化记录，流动性是给定数量、时限和价格冲击阈值下完成转换的能力。若用 $C_a(Q,\Delta t)$ 表示在时间窗 $\Delta t$ 内交易数量 $Q$ 的预期冲击成本，那么“资产 $a$ 流动”至少要同时声明 $Q$、$\Delta t$、可接受成本和市场状态。信息熵可描述订单方向或主体信念的分散程度，交易成本以货币计量，服务器与交易设施的焦耳能耗是另一类观测量；三者不能因为都被口语称作“消耗”就相互替换。由此，交易的最小表示是一个原子化状态转移
\[
(b_i,q_i,b_j,q_j,\ell_k)\xrightarrow{e_k}
(b_i',q_i',b_j',q_j',\ell_{k+1}),
\]
其中事件 $e_k$ 包含资产身份、数量、价格、双方权限和时间戳，$\ell_k$ 是清算账本状态。只有当预算、持仓、合规和清算约束同时满足时，转移才合法。这里没有必要诉诸波函数坍缩；从连续信念到确定订单，是受接口约束的离散提交。

市场的快速动力学可以用订单流而非物理粒子运动描述。令 $Q_a^{b}(p,t)$ 和 $Q_a^{s}(p,t)$ 分别为资产 $a$ 在价格档位 $p$ 的累计买量与卖量，单位是资产数量；令 $\lambda_a^{\pm}(p,t)$ 为单位时间内新增买卖订单强度，$\mu_a^{\pm}(p,t)$ 为撤单或成交强度。一个粗粒化模型为
\[
\frac{\partial Q_a^{\pm}}{\partial t}
=\lambda_a^{\pm}(p,t)-\mu_a^{\pm}(p,t)
+D_a^{\pm}\frac{\partial^2 Q_a^{\pm}}{\partial p^2},
\]
其中 $D_a^{\pm}$ 的量纲是价格平方除以时间，用于近似订单在相邻价位的迁移；它不是热扩散系数的本体等同。最优买价与卖价由非零队列的边界读出，成交事件再反馈到主体状态。若信息、订单类型、交易规则或参与者结构随时间变化，它们必须作为 $\lambda$、$\mu$、边界条件或图 $G_k^{\mathrm{con}}$ 的时变输入进入模型，而不能隐藏在一个恒定“市场势能”中。离散实现以采样步长 $\Delta t$ 更新队列，要求在该步长内事件聚合误差可接受；高频拥塞或跳价时必须改用逐事件模拟，否则连续近似会掩盖队列耗尽和价格间隙。

主体的决策同样不等于沿全局效用梯度爬升。主体 $i$ 只能基于局部信息集 $\mathcal I_i^k$、可行动作集 $\mathcal U_i^k$ 和风险约束选择订单。我们把策略写成条件核
\[
u_i^k\sim\pi_i(\cdot\mid x_i^k,\mathcal I_i^k,G_k^{\mathrm{con}}),
\qquad
x_i^{k+1}=F_i(x_i^k,u_i^k,e_i^k;\theta_i^k),
\]
其中 $u_i^k$ 可以是下单、撤单、等待或改变头寸，$e_i^k$ 是成交与外部消息，$\theta_i^k$ 是可学习参数。若把收益、风险、流动性和合规代价组成目标
\[
J_i=\mathbb E[R_i]-\alpha_i\operatorname{Risk}_i-eta_i\operatorname{Impact}_i-gamma_i\operatorname{Violation}_i,
\]
各项必须先归一到可比较的货币或无量纲评分，系数再承担单位转换和偏好权重。$J_i$ 上升只说明给定模型和观察窗下的局部目标改善；它不保证真实盈利、清算安全、社会福利或长期市场稳定。有限理性在这里表现为受限信息、受限计算、近似模型和有限行动集合，不需要引入“经济普朗克常数”。损失厌恶、现状偏好和从众也应作为可估计的策略非对称、切换成本与社会信号权重进入 $\pi_i$，而不能由某种几何方向不对易直接推出。

\label{subsec:micro_semantion}
局部主体通过订单簿和成交记录耦合，形成全局价格与流量，但聚合信号既压缩信息也制造反馈。设 $z_k$ 表示由收益率、杠杆、期限错配、抵押品折扣和网络集中度组成的宏观状态，市场清算后的更新可写成
\[
z_{k+1}=f(z_k,o_k,d_k,s_k)+\varepsilon_k,
\]
其中 $o_k$ 是订单与成交统计，$d_k$ 是违约和追加保证金事件，$s_k$ 是制度及政策输入，$\varepsilon_k$ 是未建模扰动。泡沫不是由“拓扑孤立子”这个名称证明的。可操作的正反馈链条应逐环辨识：价格上涨是否提高可抵押价值，抵押价值是否扩大融资，新增融资是否进入同一资产，强制平仓阈值是否在下跌时同步触发。我们可定义反馈增益矩阵 $K_k=\partial f/\partial z$；若在指定工作点附近其谱半径大于一，扰动可能被放大，但这只是局部线性结论。政策、主体退出、价格上限、破产和网络断边都会改变 $K_k$，因此不能从一次上涨外推全局必然崩盘。

危机中的相关性上升也不意味着资产空间真的降维或度量“坍缩”。给定收益窗口 $W_k$，可估计相关矩阵 $R_k$ 及有效维数
\[
d_{\mathrm{eff}}(k)=\frac{(\operatorname{tr}R_k)^2}{\operatorname{tr}(R_k^2)}.
\]
当共同因子占比增加时，$d_{\mathrm{eff}}$ 会下降，说明在该窗口和该估计器下可区分的波动方向减少。窗口长度、缺失数据、非同步交易和波动率变化都会影响这个量；它不能证明法律权利、产业结构或所有因果机制同时变成一维。类似地，“冻结”可被拆成交易量下降、买卖价差扩大、融资期限缩短和主体策略趋同等观测，不把流动性陷阱等同于玻璃相。物理术语可以帮助提出候选指标，但证据必须来自事件数据、资产负债表、网络暴露和政策干预。

\label{subsec:economic_cefe}
慢速制度结构的更新可在另一时间尺度上建模。令 $w_\tau\in\mathbb R^d$ 表示税率、资本约束、清算保证金、准入规则和基础设施容量等制度参数，$G_\tau^{\mathrm{con}}$ 表示契约与暴露网络；$\tau$ 以月、季度或法律事件计，而快速市场时刻 $t$ 以毫秒至日计。我们定义
\[
w_{\tau+1}=\Pi_{\mathcal W}\!left[w_\tau+eta_\tau
H(\bar z_\tau,\bar \ell_\tau,y_\tau,a_\tau)
\right],
\qquad
G_{\tau+1}^{\mathrm{con}}=\mathcal E(G_\tau^{\mathrm{con}},\mathcal R_\tau),
\]
其中 $\bar z_\tau$ 和 $\bar\ell_\tau$ 是慢时间窗统计，$y_\tau$ 是就业、产出、分配和系统风险等多目标读出，$a_\tau$ 是政策动作，$\Pi_{\mathcal W}$ 保证法规和权限边界，$\mathcal E$ 是经立法、签约、破产或基础设施建设触发的图编辑。该式替代所谓“经济认知爱因斯坦场方程”：制度演化可以有惯性、路径依赖和磁滞，但除非明确构造度量、曲率、坐标不变性和数据对应，否则黎曼张量只是一种比喻，不是经验定律。路径依赖可以通过切换成本、沉没投资、学习曲线和网络外部性检验；产业衰退则要分别测量资本折旧、技能流失、供应链退出和需求变化，而不能归因于抽象曲率自然平滑。

宏观调节因此是部分可观测、多目标并带时延的控制问题。政策制定者看到的是延迟且修订过的统计量 $y_{\tau-d}$，其动作经过银行、企业、家庭和市场结构传递。一个审慎的控制目标可以写成
\[
\min_{a_{0:T}}\sum_{\tau=0}^{T}
\bigl(omega_1 L_{\mathrm{employment}}+omega_2 L_{\mathrm{price}}
+\omega_3 L_{\mathrm{stability}}+\omega_4 L_{\mathrm{distribution}}
+\omega_5 C_{\mathrm{intervention}}\bigr),
\]
并受到权利、预算、工具权限和最坏情形风险约束。这里的权重是规范选择，不由数据自动给出；降低单一通胀损失也不代表总体任务成功。统一利率、定向工具、资本缓冲和流动性支持各有不同作用路径，算法化并不会消除模型误差。所谓“实时全息央行”还会引入隐私、歧视、单点故障、操纵和民主授权问题，因此不能把更细粒度观测直接等同于更优治理。合理的工程闭环必须保留数据最小化、独立审计、反事实评估、人工申诉和紧急退出机制。

最终，我们把市场视为由异质主体、制度结构和聚合信号共同形成的分布式适应系统，而不是具有单一信念和意志的超级个体。局部结构网络决定谁能与谁以何种权利交易；价格、利差和流量构成跨主体的全局分布表示；资产身份、合约绑定、清算耦合和政策读出把二者连接起来。市场可以完成信息聚合和资源协调，也会因激励错配、集中暴露、反馈放大与时延控制而失败。上述动力学给出可证伪的接口：改变保证金规则、限制一条融资通道、延迟部分信息或随机化政策试点，观察订单流、网络暴露和宏观状态是否按模型改变。只有通过这些干预，我们才能比较候选机制；共同曲线、物理类比或一次预测成功都不足以证明整体同构。

\section{城市交通作为图上守恒流、意图路由与多时间尺度控制}
\label{sec:section_9110}

城市交通同时包含物理移动、信息交换和制度治理，但三类过程的状态空间并不相同。我们把路网表示为随时间变化的有向多重图 $G_t^{\mathrm{road}}=(V_t,E_t)$：节点是路口、出入口和换乘点，边是具有长度、车道数、限速、通行方向和允许车型的路段。对每条边 $e$，密度 $\rho_e(x,t)$ 的单位是车辆每米，平均速度 $v_e(x,t)$ 的单位是米每秒，流量 $q_e=\rho_e v_e$ 的单位是车辆每秒；车辆 $j$ 的局部状态则包含位置、速度、尺寸、路线意图、控制模式和通信版本。道路几何与信号相位构成局部结构网络，跨路段的密度、速度和旅行时间场构成全局分布表示。车辆身份、地图匹配、路口转向绑定和传感器读出连接两类表示。排他体积来自车辆占用和安全距离，不需要类比泡利不相容；目的地是控制输入或任务条件，也不是物理电荷。

在路段内部，最小宏观模型从车辆数守恒开始。若入口、出口和事故造成的源汇为 $s_e(x,t)$，则
\[
\frac{\partial \rho_e}{\partial t}
+\frac{\partial q_e}{\partial x}=s_e,
\qquad
q_e=Q_e(\rho_e;\theta_e,t).
\]
$Q_e$ 是由道路结构、车辆组成、天气和驾驶策略决定的基本图，$\theta_e$ 是待估参数。守恒式是车辆计数结果，基本图则是模型假设，二者证据等级不同。当道路容量、限速或车道规则变化时，$Q_e$ 和边界条件都必须随时间改变。若采用长度为 $\Delta x$、步长为 $\Delta t$ 的元胞传输离散，发送量与接收量分别受上游可发车辆数和下游剩余容量约束，并需满足信息传播不跨越多个元胞的稳定条件。事故、信号突变和排队溢出时，简单平滑偏微分方程会丢失离散阻塞边界，应改用守恒数值格式或逐车模型。

逐车层可写成带反应时延的闭环控制。对车辆 $j$，令 $p_j$、$v_j$ 和 $a_j$ 分别为位置、速度和控制加速度，$d_j$ 为与前车的净间距，$v_j^{\star}$ 为受限期望速度，则
\[
\dot p_j=v_j,
\qquad
\dot v_j=a_j,
\qquad
a_j(t)=\pi_j(d_j(t-\delta_j),v_j(t-\delta_j),
\Delta v_j(t-\delta_j),u_j(t)).
\]
反应时延 $\delta_j$、最大制动、车辆长度和路面附着力必须进入安全分析。局部线性稳定只说明小扰动是否沿车队衰减；即使串稳，也不能保证路口网络无拥堵。自动驾驶或车路协同可改变策略 $\pi_j$ 和可见信息，却仍受传感误差、通信丢包、混合交通和责任规则限制。把每辆车视作一个具有全局目的函数的理性代理同样不充分：驾驶者可能临时改道，公交和急救车辆有不同优先级，行人与骑行者的安全不能压缩为小车平均速度。

路口把连续守恒流转换为离散路由和调度问题。设转向比例矩阵 $R_v(t)$ 给出进入路口 $v$ 的各入口流向各出口边的条件概率，信号动作 $g_v^k$ 从有限相位集合中选择，排队向量为 $n_v^k$。一个离散更新为
\[
n_v^{k+1}=\left[n_v^k+A_v^k-R_v^k
S_v(g_v^k,n_v^k)\right]_+,
\]
其中 $A_v^k$ 是到达量，$S_v$ 是在相位、冲突关系和饱和流率约束下的可服务量。控制器可以最小化加权排队、延误、溢出风险和相位切换成本，但必须把行人清空时间、公交优先、应急通道及公平约束写入可行域。单个路口目标下降不代表全网旅行时间下降，因为上游放行可能把车辆推入已饱和的下游边。控制实验需要同时报告局部队列、全网完成行程数、极端延误、碰撞代理指标和不同人群的可达性。

交通拥堵至少有三种不同机制。第一，需求超过割集容量，任何信号策略都无法在该时间窗内消化输入；第二，驾驶反应和换道使微小扰动放大，形成停止—启动波；第三，路由和控制的局部改进改变全网均衡，出现新道路反而提高旅行时间的布雷斯悖论。它们分别要求容量配置、局部稳定和博弈均衡分析，不能统称为“拓扑死锁”。真正的网格锁死可以通过等待依赖图定义：若一组路段相互等待下游空位且形成无可释放的有向环，则存在结构性死锁。此定义允许检测、复现与解除，例如限制进入阻塞盒、预留清空相位或暂时切断某些转向；而“度量崩塌”只提供形象描述，不能告诉控制器在哪个路口采取什么动作。

信息服务又引入第二层反馈。导航平台根据历史与实时速度估计边代价 $c_e^k$，为主体推荐路径；主体采纳后改变边流量，边流量再更新代价。令 $f^k$ 为路径流，$P(c^k)$ 为带个体偏好和合规约束的路径分配算子，则
\[
f^{k+1}=(1-\alpha)f^k+\alpha P(c^k),
\qquad
c^{k+1}=C(f^{k+1},G_k^{\mathrm{road}},w_k).
\]
$\alpha\in(0,1]$ 是更新比例，$w_k$ 是天气、活动和事故输入。过快同步改道可能造成振荡，过慢更新则保留陈旧路线。平台若只优化其用户的预计时间，可能把噪声、污染和风险转移给未使用平台的社区。由此，全局分布表示并非天然公共利益；代价函数、覆盖范围和权限边界是治理的一部分。验证导航算法时应设置未干预对照、不同采纳率、通信延迟、需求突变和地图错误，并区分估计准确、个体到达、全网效率和社会分配四类结果。

城市级控制位于更慢时间尺度。信号配时以秒或分钟更新，车道与收费规则以小时至月更新，公交线路和道路建设则以年为尺度。设快状态为 $x_k$，结构与政策参数为 $\theta_\tau$，可写成
\[
x_{k+1}=F(x_k,u_k,d_k;\theta_\tau),
\qquad
\theta_{\tau+1}=\Pi_{\Theta}
\bigl[\theta_\tau+\eta H(\bar x_\tau,y_\tau,a_\tau)\bigr].
\]
快速动作 $u_k$ 包括信号和诱导，慢速动作 $a_\tau$ 包括公交供给、拥堵收费、路权重分配与基础设施编辑。若快慢尺度没有充分分离，例如大型活动期间临时封路，慢参数必须直接进入快更新，不能沿用准稳态假设。规划目标还应包含安全、可达性、排放、残障通行、应急韧性和隐私，而不仅是平均车速。平均速度提高可能来自低速出行者被挤出系统，所以目标函数下降必须与任务成功和权利约束分别核验。

对于AgentOS工程实例，可以把路口传感摘要映射为 $h(t)$，把路线候选映射为 $E$，把信号与事件窗口映射为 $\Theta$，再由 $\Psi$ 组织跨路口读出，SPC维护短时队列上下文，TCE评估动作后果，CCM做冲突仲裁，Boundary实施安全与权限门控，Offloading调用仿真器，外部记忆 $M_e$ 保存长期流量模式。但这只是一个可执行架构：现实交通还可以采用分布式压力控制、模型预测控制、规则系统或多智能体强化学习。模块名称相同不说明城市与大脑机制相同，也不说明该实现是智能交通的唯一终局。可接受的证据来自消融、故障注入和现场对照：移除跨路口通信是否只损害协调而保留局部安全，延迟地图是否触发可预测的振荡，Boundary是否在错误目标下仍阻止危险相位。

因此，城市交通的“智能”应作操作性理解：系统利用观测和通信，在图结构与物理约束内调整路由、调度和长期结构，以改善一组公开的服务指标。它没有必要拥有统一自我；宏观有序也可能只是规则、需求和反馈共同作用的结果。局部结构网络提供道路、冲突和权限，密度与旅行时间提供全局分布表示，车辆身份、转向绑定、路口耦合和传感读出把二者联合。任何从鸟群、人脑或流体获得的类比都必须止于提出模型，随后由计数守恒、时延稳定、图割容量、现场干预和公平指标接受检验。

\section{现有语言模型：冻结参数上的条件状态传播与证据边界}
\label{sec:section_4389}

我们把这里的研究对象限定为一个无检索、无工具调用、推理期间不更新参数的自回归Transformer。这个限定排除了持续感知的具身系统，也排除了带外部记忆和验证器的复合Agent，因此结论不能直接外推到所有语言智能。设词表大小为 $V$，上下文长度为 $n$，残差维数为 $d$，层数为 $L$，参数集合为 $\theta$。输入字符串先被分词器映射成离散标识序列 $x_{1:n}\in\{1,\ldots,V\}^n$；嵌入和位置机制产生初始状态矩阵 $H^0\in\mathbb R^{n\times d}$。第 $\ell$ 层更新写成
\[
\widetilde H^\ell=H^\ell+operatorname{Attn}_\ell
(\operatorname{Norm}_\ell(H^\ell)),
\qquad
H^{\ell+1}=\widetilde H^\ell+operatorname{FFN}_\ell
(\operatorname{Norm}'_\ell(\widetilde H^\ell)).
\]
这是一个有限深度的离散复合映射，不是连续时间方程。层指标没有秒、毫秒或任何物理时间单位；只有在给出共享参数或平滑插值、步长缩放和一致极限时，才能把深度近似为连续变量。即使构造了神经常微分方程极限，也只能得到特定参数化下的连续模型，不能由形式相似推出薛定谔方程、狄拉克方程或HSF--HD整体同构。

残差状态同时编码词义、位置、句法角色、来源痕迹和任务条件，但“混合表示”并不自动等于病理。分布式表示允许一个方向参与多种读出，也允许同一功能由多个方向联合实现。我们不预设存在两个全局正交子空间，分别唯一承载“形”和“质”；这种分解若要成立，必须先指定任务、探针族和容许变换。局部结构网络可以由注意力支持、显式依存关系、实体身份及程序调用图构成；全局分布表示则体现在跨位置、跨通道的残差活动。二者通过位置身份、实体绑定、头部路由和线性或非线性读出接口联合。可解释探针若在某层读取出句法关系，只说明该信息对探针可用，不说明模型内部使用了同一坐标，也不说明该层与某个脑区机制等同。

对单个注意力头，令 $Q=HW_Q\in\mathbb R^{n\times d_k}$、$K=HW_K\in\mathbb R^{n\times d_k}$、$V_h=HW_V\in\mathbb R^{n\times d_v}$，掩码为 $M\in\{-\infty,0\}^{n\times n}$，则
\[
A=\operatorname{softmax}_{\mathrm{row}}
\left(\frac{QK^{\mathsf T}}{\sqrt{d_k}}+M\right),
\qquad
O=AV_h.
\]
$A_{ij}$ 是在该头、该层和该输入条件下，位置 $i$ 聚合位置 $j$ 的归一化权重。它可以被视为动态有向加权图，但不是黎曼度量：它通常不对称，不满足三角不等式，也会随内容改变。$A$ 也不必等于因果贡献；价值向量、输出投影、残差旁路和后续层可以放大、抵消或忽略这次聚合。把注意力称为扩散仅在额外条件下成立，例如权重非负、行随机且关注的是重复线性迭代；标准注意力还包含内容依赖、非线性归一化和多头投影，所以不能直接宣称是图拉普拉斯或狄拉克算子的离散实现。

前馈网络逐位置实施非线性变换，例如
\[
F(h)=W_2\sigma(W_1h+b_1)+b_2,
\]
它可以存储和组合统计规律，但权重不是事实数据库中的一条条可寻址记录。一个事实相关输出可能依赖嵌入、多个注意力头、多个层和最终读出；同一参数也参与许多无关任务。因而“记忆即曲率”最多是辅助解释，不能作为参数本体。模型是否掌握“巴黎—法国—首都”关系，应通过改写、反向询问、干扰上下文、知识编辑和跨语言迁移等实验检验，并报告概率变化与副作用。只在标准提示上答对，无法区分记忆、模式补全、问题泄漏或偶然采样。

推理输入并非只在初始时刻注入后消失。自回归生成在第 $s$ 步以已有上下文 $x_{1:n+s-1}$ 重新计算或增量更新键值缓存，得到下一标识条件分布
\[
p_\theta(x_{n+s}=v\mid x_{1:n+s-1})
=\operatorname{softmax}(W_Uh_{n+s-1}^L+b_U)_v.
\]
随后，解码器依据温度、截断集合、重复惩罚或约束语法选择标识。这里的温度是对logit尺度的无量纲调节，不是热力学温度；采样是随机决策提交，不是量子测量。已生成标识会成为下一步的新输入，所以系统存在内部输出回流，但在无检索设定下确实没有来自现实环境的新证据。这一区分十分关键：上下文内自反馈能够延展文本，却不能校验外部事实；采样多样性能够产生候选，却不能保证候选真实。

系统提示、用户提示和已生成历史共同构成条件输入，而不是一个在推理期间恒定的“宏观势能”。系统提示的影响会随位置、措辞、上下文竞争和模型版本变化；攻击性后缀可能覆盖前部指令，后续生成也会改变当前条件。我们可以把提示约束写成随生成步变化的控制变量 $c_s$，把缓存和残差状态写成 $z_s$：
\[
z_{s+1}=F_\theta(z_s,x_{n+s},c_s),
\qquad
x_{n+s+1}\sim D(p_\theta(\cdot\mid z_{s+1}),\omega_s),
\]
其中 $D$ 是解码规则，$\omega_s$ 是采样参数。对基础模型，$c_s$ 主要由固定上下文间接施加；对Agent，$c_s$ 还可由验证、工具结果和安全控制器更新。前者控制带宽较窄，并不等于完全没有控制；后者控制更丰富，也不自动产生正确目标或自我。

幻觉首先是输出与外部事实、给定证据或任务约束不一致的经验现象，而不是某个隐藏空间“扭曲”的同义词。它至少有五类可区分来源。其一，训练分布中的事实稀疏、冲突或过时，使条件分布没有可靠依据；其二，似然训练奖励流畅的局部续写，却未直接要求引用来源和拒答校准；其三，提示要求可能不可满足，模型仍被迫给出标识；其四，自回归早期错误进入后续上下文，形成条件级联；其五，解码与评估制度鼓励覆盖率而弱化不确定性。语义相似确实可能诱发错误关联，例如把同领域的人名、论文和年份拼接成貌似合理的引用，但不能因此断言注意力把“语义引力”误当“逻辑坐标”。要识别机制，应比较带证据和无证据提示、强制引用和自由生成、不同采样规则、逐步验证、参数编辑及反事实名称替换。

我们用三层指标诊断这类故障。标识层测量负对数似然、分布熵和采样敏感度；命题层把输出解析为带来源的断言集合，测量蕴含、矛盾、可回答性和校准；任务层检查用户目标、外部执行结果和风险约束。设输出包含命题 $y_j$，置信读出为 $r_j\in[0,1]$，事实标签为 $a_j\in\{0,1\}$，则可靠性可按分箱或Brier分数估计。低熵不等于正确，高置信也不等于有证据；相反，较高熵有时只是存在多个同义表达。只有将标识分布映射到命题身份并连接外部证据，才能讨论事实校准。模型目标下降、下一词预测正确与任务成功处于不同层级，不能互相代替。

\label{subsec:llm_spatial_topology_pathology}
Transformer的层堆叠可以视作一组具有不同参数的功能算子，但不能把每个Block直接等同于固定脑区或独立子流形。对输入 $x$，第 $\ell$ 层的有效雅可比为 $J_\ell(x)=\partial H^{\ell+1}/\partial H^\ell$；跨层扰动传播由矩阵积 $J_{L-1}\cdots J_0$ 决定。若某些奇异值长期接近零，相关方向会被压缩；若远大于一，小扰动可能放大。层归一化、残差连接和注意力会使这个谱随输入变化，因此“串行层数越多等于心理时间越长”不是普遍结论。增加深度可以增加复合计算步骤，也可能造成过平滑、梯度困难或冗余；循环使用较小模型则可能通过外部状态和多次调用获得额外计算。空间参数量与推理步数之间存在资源替代，但不存在仅凭维数就能否定小模型时间循环的定理。

所谓“拓扑死锁”也必须落到可观测故障。若模型在重复生成环、指令冲突或注意力汇处无法离开，可以定义状态重复率、有效新信息率、约束违例次数和停止条件触发率。干预包括改变位置编码、屏蔽特定头、重置缓存、插入验证标识和外部中断。只有当某种结构干预稳定地解除故障，才能把原因归于相应路由机制。把所有失败归因于静态基质，会忽略数据、解码器和提示制度；把所有成功归因于高维流形，又会失去可证伪性。

\label{subsec:llm_wm_coherence_capacity}
上下文容量不能只用窗口标识数衡量。我们定义任务相关相干容量 $C_{\varepsilon}(T,\mathcal P)$：在任务分布 $\mathcal P$、误差阈值 $\varepsilon$ 和指定提示协议下，模型仍能保持实体身份、约束关系、时间顺序和目标条件的最大有效信息量。这个量依赖分词、位置机制、干扰模式、检索方式和读出标准。可将上下文中的必要约束构成集合 $R=\{r_1,\ldots,r_m\}$，在不同距离和干扰条件下测试每个约束的保持概率 $p_j$，并用
\[
C_{\mathrm{bind}}=\sum_{j=1}^{m}w_jp_j,
\qquad \sum_jw_j=1,
\]
形成无量纲绑定得分。它测量的是任务协议下的可靠读出，不是模型脑容量或焦耳能量。长窗口扩大可访问输入 $I$，参数和计算扩大候选处理能力 $C$，但两者都可能增加竞争与校验负担，使单位输出的可靠带宽下降。

\label{subsec:llm_wm_coherence_capacity_extended}
相干容量还受运行时预算限制。设第 $s$ 步可用计算量为 $B_s$，待维护关系数为 $m_s$，平均检索歧义为 $a_s$，外部验证成本为 $v_s$。一个经验风险模型可写成
\[
\Pr(\text{binding error at }s)
=\sigma(\beta_0+\beta_1m_s+\beta_2a_s+\beta_3v_s
-\beta_4B_s+\beta_5\delta_s),
\]
其中 $\delta_s$ 是相关证据距当前读出位置的有效距离，$\sigma$ 为逻辑函数。该式是待拟合的统计模型，不是物理定律；系数会随模型和协议改变。它允许比较“扩大窗口”“压缩历史”“检索少量证据”“增加验证计算”四种策略，而不是先验宣称空间必然挤压时间。实验应固定任务难度，随机排列无关干扰，分别测量实体错绑、约束遗失、循环引用和最终答案准确率。

\label{subsec:cognitive_thermodynamics_cct_dcm}
动态相干管理可以被定义为有限预算下的状态选择，而不是热力学熵泵。令候选上下文单元集合为 $U_s$，每个单元有相关性 $r_u$、来源可靠性 $e_u$、新颖性 $n_u$、标识成本 $c_u$ 和绑定风险 $b_u$；选择变量 $g_u\in\{0,1\}$。运行时管理器求解
\[
\max_{g}\sum_{u\in U_s}g_u
(\alpha r_u+\beta e_u+\gamma n_u-\lambda b_u)
\quad\text{s.t.}\quad
\sum_{u\in U_s}g_uc_u\le B_s,
\]
并维护必须保留的身份、目标和安全约束。实际实现可用启发式、学习排序或约束解码近似，不要求每步精确求解。若目标、任务结构或预算随时间变化，评分和约束必须进入每次更新。摘要、滑动窗口、缓存淘汰、分层记忆、检索和思维链都可看作不同选择策略；它们可能降低上下文负担，也可能删除关键反例或放大错误摘要，因此必须与完整上下文、随机选择和人工上界比较。

在AgentOS实例中，$h(t)$承载当前残差、命题和工具返回的有限界面，$E$保存带时间与来源的事件，$\Theta$保存较慢的结构与策略，$\Psi$提供任务相关自我和承诺读出；SPC压缩当前状态，TCE调节探索与控制，CCM估计输入和运行复杂度，Boundary限制数据与行动权限，Offloading把可复用过程交给脚本、文件和外部记忆$M_e$。这些模块为动态相干管理提供一种明确落点，但基础Transformer、循环状态模型、检索系统和人机协同也可用别的机制实现相似功能。工程比较应基于相同任务、时延、费用和风险预算，而不是以是否采用AgentOS术语划分智能等级。

\label{subsec:tda_manifold_lithography}
拓扑数据分析可以探测隐藏状态点云在选定度量和尺度下的形状，但其结论具有表示依赖性。给定样本集合 $X=\{x_i\}$、层 $\ell$ 和读出位置 $p$，收集向量 $h_i^{\ell,p}$，先固定中心化、归一化和距离 $d$，再按半径 $\epsilon$ 构造Vietoris--Rips复形。持续同调追踪连通分支、环和高维空腔在尺度变化中的出生与死亡。长寿命条形说明该采样和度量下存在稳定拓扑特征；它不自动对应概念、逻辑漏洞或意识结构。样本分布、维数压缩、各向异性和有限数据都会制造或抹去孔洞，所以报告必须包含重采样置信区间、随机标签基线、距离选择和层间对照。

Attention Sink也不应直接称为虫洞。我们可操作地定义某些位置在多层多头上获得异常高的入边权重，随后检验屏蔽、替换或移动这些位置是否改变信息保持和任务结果。如果汇聚权重高但因果干预影响小，它只是可视化显著而非功能枢纽；如果干预同时破坏有用路由和错误生成，则需区分容量管理与故障来源。过度平滑可通过位置间表示方差下降、类别可分性下降和不同输入输出趋同来测量，但残差与多头机制可能保留局部差异。任何“拓扑奇点”诊断都应给出替代统计量，避免把漂亮的二维图当作机制证明。

\label{subsec:tda_manifold_evolution}
以真值判断研究为例，若在中间层观察到真假命题可被低维方向分离，这是一项关于特定数据集和线性探针的经验结果。它说明模型状态包含对任务有用的信号，不说明真理本身位于一个固定几何轴，也不说明模型会在生成时使用该轴。更严格的协议应训练于一组主题、测试于未见主题，交换表述和标签，控制句长、否定词与来源，并用方向消融或激活编辑检验因果作用。若探针准确而编辑无效，信号可能只是旁观读出；若编辑提高一组任务却损害另一组任务，则该方向不是无条件的“真值控制杆”。

Mapper等方法可把点云划分成重叠区间并构造簇图，用于发现候选状态转移。滤波函数、覆盖宽度、重叠率和聚类算法都会改变图形，因此必须做参数敏感性分析。观察到从多簇到少簇的层间变化，可以解释为任务相关压缩、类别聚合或过平滑的候选证据；仅凭图像不能区分这些机制。对幻觉样本，我们应比较正确与错误生成的路径、来源证据强度、校准和解码随机性，并在相同命题身份上进行配对，而不是把错误点落入某个孔洞就宣布捕获了幻觉本体。

\label{subsec:global_manifold_mapping}
全量测绘在工程上还面临样本覆盖与计算规模。模型可能在不同语言、任务、位置和生成阶段形成互相重叠的局部坐标，单一二维嵌入会压扁这些条件。更稳妥的做法是建立分层图谱：底层保存输入身份、层、位置、头和模型版本；中层保存局部邻域、探针与干预结果；上层才汇总可重复的功能区和转移模式。跨版本比较需要对齐共同任务和读出，而不是强行逐神经元匹配。所谓“多孔拓扑”只有在多次采样、多个尺度和外部任务上复现，且干预孔洞边界会改变能力或故障时，才获得机制意义。

由此，我们对冻结语言模型作出有限结论。它是一个在固定参数上运行的条件概率计算系统，能够借助分布式状态和动态路由完成丰富的模式组合；它也受证据缺失、自回归级联、绑定容量和运行预算约束。参数冻结不等于系统完全静止，隐状态和缓存会变化；隐状态变化也不等于结构学习，参数与持久外部记忆并未因此更新。现有LLM可以作为HSF--HD一般状态动力学的一个对象，用状态、结构、读出、目标和边界进行分析，但预测有效不能推出整体同构。只有加入可追踪证据、反事实验证、运行时复杂度管理和受治理的外部行动后，它才构成更完整的工程闭环；这仍是能力组合，不是关于意识或唯一AGI架构的证明。
